\section{Implementation}\label{sec:impl}
In this section, we describe the implementation detail of \( \tool \),
which is an open-source project\footnote{The URL of the tool is anonymized
due to a double-blind review process.}.

Our goal is to extract the JavaScript syntax and semantics from the
ECMAScript 2020 specification.  It is the next ECMAScript version
that is actively updated in its GitHub repository~\cite{es2020}
to be released in the next year.  Moreover, Ecma Technical Committee 39 (TC39)
officially provides the conformance test suite, test262, for ECMAScript 2020.

\subsection{Excluded Language Features}
We decided not to support the following language features in ECMAScript 2020:

\paragraph{No-Strict Mode}
Since ECMAScript 5, JavaScript introduced strict mode to
\textit{intentionally} provide slightly different semantics.  It
provides more stable semantics by preventing unsafe actions or
throwing originally ignored errors.  However, the ECMAScript
specification states that ``Strict mode does not restrict or modify
any aspect of the ECMAScript semantics that must operate consistently
across multiple source text units.''  Thus, we decided to support only
strict mode of JavaScript to focus on the main contribution of this
paper.

\paragraph{Orthogonal Built-in Objects}
JavaScript provides various standard built-in objects to support
helper functions.  However, several built-in libraries are orthogonal
to the JavaScript semantics but support extra functionalities such as
\( \code{Math} \), \( \code{Date} \), \( \code{RegExp} \), and
\( \code{JSON} \).  Thus, we decided not to support the four libraries.

\paragraph{Modules}
Moreover, we do not support modules that were introduced since ES6
(ECMAScript 2015).  With the \( \code{import} \) keyword, a JavaScript
program can import other JavaScript components defined with the \(
\code{export} \) keyword in their modules.  In order to support
modules in our automatically generated semantics, we should model the
functionality to communicate with file systems to access JavaScript
module files.  Because modules do not affect the overall JavaScript
semantics critically, we decided not to support modules.

\subsection{Compile Rules for ECMAScript}

\begin{table}[t]
  \centering
  \caption{Compile rules in ECMAScript 2020}
  \label{table:rules}
\vspace*{-1em}
  \[
    \begin{array}{c|r|r|r|r|r|r|r}
      \text{name}
      & \multicolumn{1}{c|}{\code{Stmt}}
      & \multicolumn{1}{c|}{\code{Expr}}
      & \multicolumn{1}{c|}{\code{Cond}}
      & \multicolumn{1}{c|}{\code{Value}}
      & \multicolumn{1}{c|}{\code{Ty}}
      & \multicolumn{1}{c|}{\code{Ref}}
      & \multicolumn{1}{c}{\code{SecNo}}\\\hline
      \text{\# rules}
      & \inred{XX}
      & \inred{XX}
      & \inred{XX}
      & \inred{XX}
      & \inred{XX}
      & \inred{XX}
      & \inred{XX}\\\hline
    \end{array}
  \]
\end{table}

As we discussed in Section~\ref{sec:compiler}, \textsf{Algorithm Compiler}
requires compile rules to compile given abstract algorithms to \( \ires \)
functions.  In order to cover the entire ECMAScript 2020 specification,
we defined compile rules for only seven kinds as summarized in
Table~\ref{table:rules}.  The compile rule for statements, \( \code{Stmt} \),
contains \inred{XX} compile rules and generates \( \ires \) instructions.
The \( \code{Expr} \), \( \code{Cond} \), and \( \code{Value} \) compile
rules denote expressions, conditions, and values, respectively.  While
all of them generate \( \ires \) expressions, they represent different
contexts in the ECMAScript specification.  For example,
\( \code{Expr} \) represents a context where any expression can appear,
\( \code{Cond} \) denotes a context where any boolean-valued expression
can appear, and \( \code{Value} \) represents a context where a fully
evaluated value can appear.  The \( \code{Ty} \) compile rule denotes
type names and generates string primitives in \( \ires \).  The \( \code{Ref} \)
represents references such as identifier lookup and member accesses of
map-like structures.  It generates \( \ires \) reference expressions.
Finally, the \( \code{SecNo} \) compile rule denotes section numbers
like \( \code{9.5.12} \), and it converts such numbers to
corresponding name strings like \( \code{"ProxyExoticObject.Call"} \):
\begin{lstlisting}[style=myScalastyle]
lazy val SecNo =
  "9.5.12" ^^ { case x => "ProxyExoticObject.Call" } |||
  ...
\end{lstlisting}

\subsection{Global Setting}
We manually implemented \textsf{Global Setting} for ECMAScript 2020,
which has two parts: ECMAScript data types and built-in objects.

\paragraph{ECMAScript Data Types}
In ECMAScript 2020, Section 6 \textsf{ECMAScript Data Types and Values}
describes data types used in the specification.  It describes not only
data types explicitly exposed for JavaScript evaluation, but also all
the internal data types used in the specification.  For example, the
Environment Records are record types from identifiers to bindings for
values.  They are not explicitly accessible in JavaScript programs
but defined and used internally in the specification.

Each ECMAScript data type has several fields to store their relevant values.
We should implement such structures in order to correctly evaluate JavaScript
programs, but fields of ECMAScript data types could be abstract algorithms.
We call them \textit{method-like abstract algorithms} because they
implicitly get their receiver objects as arguments at callsites.
To mimic such implicit behaviors, we added a special variable
\( \code{this} \) as the first parameter of a method-like abstract algorithm,
and passed a receiver object at its callsite by modifying
\textsf{Algorithm Compiler}.  For example, an Environment Record has
the \textsf{DeleteBinding (N)} abstract algorithm as one of its fields.
Then, it gets two parameters, the special parameter \( \code{this} \)
and normal one \( \code{N} \).  Thus, the following algorithm step:
\begin{center}
  \includegraphics[width=0.22\textwidth]{img/method-like-example.png}
\end{center}
is compiled into the following \( \ires \) instruction:
\begin{lstlisting}[style=ires]
return DclRec.DeleteBinding(DclRec, N)
\end{lstlisting}

\vspace*{-.5em}
\paragraph{Built-in Objects}
In ECMAScript, built-in objects are pre-defined functions with several
built-in functions.  For example, \( \code{Array} \)
is the constructor of array objects, and its prototype
\( \code{Array.prototype} \) has built-in functions for array objects.
For instance, \( \code{[1,2,3].flat()} \) calls the
\( \code{Array.prototype.flat} \) built-in function with the array
\( \code{[1,2,3]} \).  Because built-in functions are also abstract
algorithms, each of them is automatically converted to an \( \ires \)
function.  However, the structures of built-in objects should be
manually implemented.  Thus, we implemented built-in objects in Scala
and connected their properties with the extracted \( \ires \) functions.

Some built-in objects that are explicitly referenced in abstract
algorithms are intrinsic objects.  Intrinsic objects have their own
aliased names summarized in Table 7 of Section 6.1.7.4
\textsf{Well-Known Intrinsic Objects}.  We extracted the alias
information into \textsf{Global Setting} to utilize it during evaluation.
